\documentclass[11pt]{article}
\usepackage[margin=2.2cm]{geometry}
\usepackage{booktabs}
\usepackage{array}
\usepackage{amsmath}
\usepackage{xcolor}
\usepackage{sectsty}
\usepackage{enumitem}
\usepackage{tcolorbox}
\usepackage{fancyhdr}
\usepackage{titlesec}
\usepackage{longtable}
\usepackage[hidelinks]{hyperref}
\definecolor{accent}{RGB}{40,90,140}
\definecolor{warn}{RGB}{150,60,30}
\definecolor{soft}{RGB}{95,100,110}
\allsectionsfont{\color{accent}\sffamily}
\titleformat{\section}{\Large\bfseries\sffamily\color{accent}}{\thesection}{0.6em}{}
\titleformat{\subsection}{\large\bfseries\sffamily\color{accent}}{\thesubsection}{0.5em}{}
\renewcommand{\arraystretch}{1.25}
\pagestyle{fancy}\fancyhf{}
\rhead{\sffamily\small\color{soft}ufpipe — Ulsan-Fault catalog pipeline reference}
\rfoot{\sffamily\small\color{soft}\thepage}
\setlength{\headheight}{14pt}
\newtcolorbox{keybox}[1]{colback=accent!6,colframe=accent!55,title=#1,fonttitle=\bfseries\sffamily,
  left=3mm,right=3mm,top=1.5mm,bottom=1.5mm}
\newtcolorbox{warnbox}[1]{colback=warn!6,colframe=warn!55,title=#1,fonttitle=\bfseries\sffamily,
  left=3mm,right=3mm,top=1.5mm,bottom=1.5mm}
\newcommand{\code}[1]{\texttt{#1}}
% let long typewriter paths break inside words so they don't overflow the right margin
\usepackage{url}
\sloppy
\emergencystretch=3em

\begin{document}

\begin{center}
{\huge\bfseries\sffamily\color{accent} Ulsan-Fault earthquake catalog pipeline}\\[2mm]
{\large\sffamily Reference manual: detection $\rightarrow$ association $\rightarrow$ location $\rightarrow$ relocation (HypoDD dt.ct + dt.cc)}\\[2mm]
{\small\color{soft} Long-term (2010–present) catalog for the Ulsan Fault, SE Korea\hfill rev.\ 2026-07-27 · \code{src/ufpipe}}
\end{center}
\vspace{1mm}\hrule\vspace{3mm}

{\small \code{ufpipe} (\code{src/ufpipe}) is the \textbf{end-to-end pipeline}: six stages, detection $\rightarrow$
association $\rightarrow$ augment $\rightarrow$ PHS $\rightarrow$ locate $\rightarrow$ \textbf{relocate} (HypoDD
dt.ct $+$ dt.cc). The final \code{relocate} stage produces the double-difference relocated catalog. The heavy
dt.cc/xcorr/HypoDD \emph{engine} is external (\code{15.PocketQuake} korea-cluster-relocation, HypoDD v2.1beta);
ufpipe orchestrates it. The relocate stage is \textbf{self-fed}: it builds its inputs from ufpipe's OWN
detection+association (\code{ufpipe.reloc\_inputs}) and runs the orchestration in
\code{src/ufpipe/reloc\_driver/}; working dirs + results go to \code{outputs/reloc/}.
\code{detection\_test/} is a \textbf{fully frozen pilot archive} (the 2016 4-picker comparison;
\code{lib/} DEPRECATED, see \code{detection\_test/lib/DEPRECATED.md}). The relocation internals have a
dedicated guide: \code{src/ufpipe/reloc\_driver/PIPELINE.md} + the pilot's \code{study\_guide.pdf}.}

\begin{center}\small\color{soft}
\textbf{Data scope (verified 2026-07):} the waveform archive covers \textbf{2010--2024}, ending at
2024 day 305 (31 Oct 2024). No 2025+ data is present in any network directory, so 2024's 305 daily pick
files are complete for the data that exists.
\end{center}

\tableofcontents
\vspace{3mm}\hrule

%================================================================================
\section{What the pipeline does}

Build an \textbf{internally consistent, long-term earthquake catalog} for the Ulsan Fault region (SE Korea),
2010–present, from continuous seismic waveforms, using AI phase pickers so small events are detected uniformly
even as the recording network grows, and finally \textbf{relocate} the catalog with double-difference (HypoDD
dt.ct + dt.cc) for sub-100\,m relative precision. Every stage is scripted and parameterized; results reproducible.

\begin{center}
\small
\begin{tabular}{@{}ccccc@{}}
\textbf{Detection} & \textbf{Association} & \textbf{Location} & \textbf{QC} & \textbf{Relocation} \\
PhaseNet / PhaseNet+ & PyOcto (+augment) & HYPOINVERSE & \code{uf\_cluster} & HypoDD dt.ct + dt.cc \\
\emph{picks} $\rightarrow$ & \emph{events+assign} $\rightarrow$ & \emph{.sum/.arc} $\rightarrow$ & \emph{well-located} $\rightarrow$ & \emph{hypoDD.reloc} \\
\end{tabular}
\end{center}

\subsection{\code{ufpipe} vs \code{detection\_test}}
\code{ufpipe} is the production end-to-end pipeline; everything live is under \code{src/}, including the
relocation orchestrator \code{src/ufpipe/reloc\_driver/run\_picker\_reloc.py} (driven by the relocate stage
with \code{-{}-skip-build}; the HypoDD/xcorr engine itself is external in \code{15.PocketQuake}).
\code{detection\_test} is the \textbf{frozen pilot archive} of the 2016 4-picker comparison — nothing in it
is live (\code{lib/} DEPRECATED; the \code{reloc\_2016\_uf*/} dirs hold the pilot's manifests and results).
\begin{center}
\small
\begin{tabular}{@{}>{\raggedright\arraybackslash}p{2.9cm}>{\raggedright\arraybackslash}p{5.6cm}>{\raggedright\arraybackslash}p{5.7cm}@{}}
\toprule
 & \textbf{\code{ufpipe}} (\code{src/ufpipe}) & \textbf{\code{detection\_test}} \\
\midrule
Role & the self-contained end-to-end catalog + relocation pipeline (incl.\ \code{reloc\_driver/}) & frozen archive of the 2016 4-picker pilot \\
Stages & detection $\rightarrow$ association $\rightarrow$ augment $\rightarrow$ PHS $\rightarrow$ locate $\rightarrow$ \textbf{relocate} & --- (nothing live) \\
Driver & \code{python -m ufpipe.run\_pipeline} & --- \\
Ends at & HypoDD dt.cc (\S7); results in \code{outputs/reloc/} & --- \\
\bottomrule
\end{tabular}
\end{center}
ufpipe's own detection + association cover \textbf{KS/KG/GJ/NS} with daily-chunked association (\S5), and the
\code{relocate} stage builds its inputs from those outputs (\code{ufpipe.reloc\_inputs}); it no longer consumes
\code{detection\_test/lib}. The pipeline is fully self-contained. See \S7.

\subsection{Two orthogonal ``model'' dimensions}
\begin{center}
\small
\begin{tabular}{@{}lll@{}}
\toprule
Dimension & Flag & Values \\
\midrule
\textbf{Picker model} & \code{-{}-model} & \code{stead}, \code{original}, \code{phasenet\_plus}, \code{eqt}, \dots \\
\textbf{Velocity model} & \code{-{}-velmodel} & \code{kim1983}, \code{kim2011} \\
\bottomrule
\end{tabular}
\end{center}
Any picker model can be located with any velocity model. \code{stead}/\code{original} are SeisBench PhaseNet
weights; \code{phasenet\_plus} is EQNet PhaseNet+ (needs a local EQNet clone).

%================================================================================
\section{Install \& environment}

The pipeline runs in a \textbf{two-env split} on this server (there is no single ``ulsan'' env):
detection (PhaseNet+/SeisBench, torch) in \code{eqnet} (Py 3.9); association (PyOcto) + orchestration in
\code{base} (Py 3.12); the relocate stage's GPU xcorr shells out to \code{pq-gpu} internally.

\begin{keybox}{One-time setup}
\begin{enumerate}[leftmargin=6mm,itemsep=0.5mm]
\item \code{conda run -n eqnet pip install -e . -{}-no-deps} \hfill(detection env)
\item \code{conda run -n base~ pip install -e . -{}-no-deps} \hfill(association/orchestration env)
\item \textbf{HYPOINVERSE}: the external \code{hyp1.40} binary must be on \code{PATH} (not pip-installable).
\item \textbf{PhaseNet+ only}: a local EQNet clone at \code{config.EQNET\_DIR}.
\item \textbf{Fresh machine}: \code{conda env create -f environment.yml} (env \code{uf-catalog}), then install the
      torch build matching your CUDA (\code{pip install torch -{}-index-url https://download.pytorch.org/whl/cu128}).
\end{enumerate}
\end{keybox}

After \code{pip install -e .}, \code{import ufpipe.config} and \code{python -m ufpipe.run\_pipeline} work from any
directory — no \code{sys.path} hacks.

%================================================================================
\section{Stages 1--5: detection through location}

\code{STAGES = [detection, association, augment, phs, locate, relocate]} — run in order. Stages 1--5 (below) are
functions in \code{src/ufpipe/core.py}; the relocate stage (6) is \code{src/ufpipe/relocate.py} and has its own
\S7. \code{run\_pipeline.py} chains them per year.

\begin{center}
\small
\begin{longtable}{@{}c>{\raggedright\arraybackslash}p{2.4cm}>{\raggedright\arraybackslash}p{5.0cm}>{\raggedright\arraybackslash}p{5.2cm}@{}}
\toprule
\# & Stage & In $\rightarrow$ Out & What it does \\
\midrule
\endhead
1 & \textbf{detection} & continuous mSEED $\rightarrow$ \code{picks/} (per station-day) & AI picker on each station-day across \textbf{KS/KG/GJ/NS} (per-year multi-network table, \code{stations.py}); anti-alias $>$100 Hz, unify mixed rates, resample to 100 Hz, pick P/S above threshold. GPU. Idempotent (skips done days). \\
2 & \textbf{association} & picks $\rightarrow$ \code{pyocto\_*.csv} + assignment & PyOcto octree associator groups picks into events (gate below), \textbf{daily-chunked} ($\pm$150 s/day, in-day origins) so it scales to the dense NS array. Coords from the multi-network table. \\
3 & \textbf{augment} & assignment $\rightarrow$ augmented assignment & scan daily picks for orphans PyOcto missed, add to each event so HYPOINVERSE sees the full pick set. Backs up + overwrites in place; idempotent. \\
4 & \textbf{phs} & assignment $\rightarrow$ \code{UF\{year\}.phs} & write the HYPOINVERSE phase file from the (augmented) assignment. \\
5 & \textbf{locate} & \code{.phs} $\rightarrow$ \code{UF\{year\}.\{sum,arc\}} & run HYPOINVERSE (\code{hyp1.40}) with the chosen velocity model $\rightarrow$ located catalog + errors. \\
\bottomrule
\end{longtable}
\end{center}

\begin{warnbox}{Pipeline order matters}
\code{augment} runs \emph{between} \code{association} and \code{phs}, so HYPOINVERSE sees the full augmented pick
set and computes correct \code{GAP}/\code{DMIN}/\code{ERZ}/\code{RMS}. Do not reorder.
\end{warnbox}

\subsection{Source of truth for pick$\leftrightarrow$event assignment}
The \textbf{PyOcto assignment (after augmentation)} is authoritative for which \code{(station, phase)} tuples
belong to each event. Do NOT use a legacy time-window pick dump — when two events fall $<$60\,s apart at the same
station the earliest pick ``wins'' and leaks into the wrong event. Downstream (SAC-header writers, magnitude/QC
notebooks) must read \code{pyocto\_assignment\_\{velmodel\}\_\{year\}.csv}.

%================================================================================
\section{Parameters (\code{src/ufpipe/config.py} — the single disclosed source)}

Change defaults here. Everything is one file.

\subsection{Detection / DSP}
\begin{center}
\small
\begin{tabular}{@{}>{\raggedright\arraybackslash}p{3.6cm}>{\raggedright\arraybackslash}p{4.0cm}>{\raggedright\arraybackslash}p{6.6cm}@{}}
\toprule
Parameter & Value & Meaning \\
\midrule
\code{SAMPLING\_RATE} & 100.0 Hz & the pickers' trained rate; all data resampled to it \\
\code{ANTIALIAS} & lowpass 45 Hz, 8 corners, zerophase & applied BEFORE down-sampling ($>$100 Hz KG/NS) \\
\code{BANDPASS} & 1--40 Hz, 4 corners, causal & legacy detection path; PhaseNet+ runs raw \\
\code{P\_THRESHOLD}, \code{S\_THRESHOLD} & 0.2 & SeisBench pick probability threshold \\
\code{PNPLUS\_MIN\_PROB} & 0.3 & EQNet PhaseNet+ pick threshold \\
\code{PNPLUS\_NT} & $1024{\times}36$ & PhaseNet+ patch length (whole-day, \code{cut\_patch}) \\
\code{MAX\_CORES} & 24 & CPU cap (respects \code{taskset} affinity) \\
\bottomrule
\end{tabular}
\end{center}

\subsection{Association (PyOcto) — daily-chunked, \code{ASSOC\_GATE} (lenient) vs \code{ASSOC\_GATE\_STRICT}}
Association runs \textbf{one calendar day at a time}: a $\pm$\code{ASSOC\_OVERLAP\_S} (150 s) window per day is
associated and only events whose \emph{origin} falls in-day are kept (dedup). This is physically equivalent to a
single pass (local events are seconds long) but keeps each PyOcto solve to seconds --- a whole-year single pass is
intractable on the dense $\sim$200-station NS array ($\gg$1 h, 12 GB). Station coordinates come from the
multi-network year table (\code{stations.py}), so KS/KG/GJ/NS all associate.
\begin{center}
\small
\begin{tabular}{@{}lcc l@{}}
\toprule
Parameter (\code{config}) & lenient & \code{\_STRICT} & Meaning \\
\midrule
\code{REGION\_CENTER} $\pm$ pad & \multicolumn{2}{c}{(35.856,129.224) $\pm$ (1.0,1.2)} & search area (deg) \\
\code{ASSOC\_ZLIM} & \multicolumn{2}{c}{(0, 30)} & depth range (km) \\
\code{ASSOC\_TIME\_BEFORE} & \multicolumn{2}{c}{300} & origin-search window (s) \\
\code{ASSOC\_OVERLAP\_S} & \multicolumn{2}{c}{150} & daily-chunk overlap (s) \\
\code{n\_picks} & 4 & 6 & min total picks \\
\code{n\_p} / \code{n\_s} & 2 / 2 & 3 / 3 & min P / S picks \\
\code{n\_ps} & 1 & 2 & min stations with both P+S (depth) \\
\code{ASSOC\_PICK\_MATCH\_TOL} & \multicolumn{2}{c}{1.5} & residual cap (s) \\
\bottomrule
\end{tabular}
\end{center}
Every value above is \textbf{overridable per run} (defaults = config) with matching flags on
\code{association.py} and \code{run\_pipeline.py}: \code{-{}-gate n\_picks=…,n\_p=…,n\_s=…,n\_ps=…},
\code{-{}-pick-match-tol}, \code{-{}-overlap-s}, \code{-{}-zlim Z0,Z1}, \code{-{}-center LAT,LON},
\code{-{}-lat-pad}, \code{-{}-lon-pad}, \code{-{}-time-before}; \code{-{}-strict} selects
\code{ASSOC\_GATE\_STRICT} (ignored if \code{-{}-gate} is given). The cockpit notebook sets the same knobs via
\code{run\_association\_year(overrides=…)}. Velocity = \code{ASSOC\_VELMODEL} (\code{kim2011} 1-D).
\textbf{Station-code alias}: a picked code missing from the year table but present as \code{<code>2}
(e.g.\ \code{BUS}$\to$\code{BUS2} — same station, renamed code) is remapped so its picks are used, not dropped.
Output schema (events \code{idx,time,x,y,z,picks,latitude,longitude,depth}; assignment
\code{event\_idx,pick\_idx,residual,station,phase,time}) is unchanged, so \code{phs}/\code{locate}/\code{relocate}
consume it as before.

\smallskip\noindent\textbf{HYPOINVERSE pick weights (stage 4 input).} \code{make\_phs} sets the HYPO71 weight
code from the picker's own probability (\code{config.PHS\_WEIGHT\_SCHEME="probability"}): \code{PHS\_WEIGHT\_BINS}
$\ge$0.90$\to$0, $\ge$0.70$\to$1, $\ge$0.50$\to$2, $\ge$0.30$\to$3, else 4 (0 = full weight); S picks one code
worse (\code{PHS\_WEIGHT\_S\_PENALTY}). This matches the relocation stage's \code{phs\_weight\_scheme="probability"}.
Set \code{PHS\_WEIGHT\_SCHEME="phase"} for the legacy fixed P=0/S=1.

\subsection{Velocity model \& location}
\code{DEFAULT\_VELMODEL = kim2011}; \code{kim1983} also available (\code{data/metadata/velocity/kim1983.csv},
a layered $V_p$/$V_s$ model; HYPOINVERSE \code{*.crh} tables live in \code{data/hypoinv/\{kim1983,kim2011\}/}).
\code{PHASE\_CHANNELS = \{P: HHZ, S: HHN\}}.

\medskip
\noindent\textbf{The HYPOINVERSE station file is GENERATED, not hand-maintained.} Before locating,
\code{stations.write\_hypoinverse\_sta(year)} derives \code{STA/UF<year>\_hyp.sta} (fixed-format) and
\code{STA/UF<year>.sta} (CSV, read by \code{uf\_cluster.load\_stations}) from the multi-network year table ---
the single source of truth. \code{-{}-no-regen-sta} keeps an existing file; the previous hand file is
preserved once as \code{*.legacy}.

\begin{warnbox}{Why this matters --- a silent and destructive failure mode}
A stale or incomplete station list does \emph{not} raise an error. HYPOINVERSE prints
\code{SKIP PHASE CARD WITH UNKNOWN STATION} for every pick at a station it does not know, then drops any event
left with $<4$ usable phases (\code{CANT SOLVE}) --- and those losses propagate into the \code{.arc}
$\rightarrow$ ph2dt $\rightarrow$ \code{event.dat}/\code{dt.ct} that HypoDD relocates on. The shipped legacy
list held \textbf{10} stations for 2010 while the association used \textbf{12}: 21\% of picks were discarded.
Regenerating took 2010 from \textbf{640 located / 161 QC-pass} to \textbf{791 / 303}. The gap was far larger in
later years --- the legacy list knew 13 of 46 stations in 2016 (no GJ array) and 44 of 243 in 2021 (no NS array).
\end{warnbox}

\subsection{Instrument responses --- all four networks (\code{src/ufpipe/responses.py})}
\code{responses.py} is to responses what \code{stations.py} is to coordinates: the single source of truth.
\code{load\_inventory(networks=None)} returns one merged \code{obspy.Inventory} spanning KS/KG/GJ/NS, built
master-first so an authoritative entry is never shadowed by a per-station fill. Sources, all under
\code{data/metadata/responses/}:

\begin{center}\small
\begin{tabular}{@{}llll@{}}
\toprule
\textbf{Network} & \textbf{Directory} & \textbf{Format} & \textbf{Coverage} \\
\midrule
KS + KG & \code{master/} & StationXML & \code{KS\_KG\_metadata\_1.0.2.xml} \\
KS (7 gaps) & \code{fetched/extracted/} & SEED RESP & fetched from NECIS \\
GJ & \code{gj/} & SEED RESP & 30 stations $\times$ 3 (MARA) \\
NS & \code{ns/} & SEED RESP & 200 stations $\times$ 3 (MARA) \\
NS \code{N201}--\code{N220} & \code{ns\_derived/} & SEED RESP & 20 stations $\times$ 3, \textbf{derived} \\
\bottomrule
\end{tabular}
\end{center}

\noindent Coverage is \textbf{100\% of the station table for every year 2010--2024}, all three channels.
Check it with \code{python -m ufpipe.responses -{}-coverage <year>}. The ML side reaches the same inventory
through \code{ml\_pipeline.load\_full\_inventory()}.

\begin{warnbox}{A missing response is silent rather than loud}
Response removal is what turns counts into Wood-Anderson millimetres. A station absent from the inventory
raises no error --- \code{remove\_response} simply skips it, so the event's ML is the median over whatever
stations happened to have metadata. Before GJ and NS responses existed, an ML computed on a 2016 or 2021
catalogue would have quietly used only the KS/KG minority (7 of 46 stations in 2016; 43 of 243 in 2021)
while appearing to succeed. Run the coverage check before trusting any ML catalogue.
\end{warnbox}

\begin{warnbox}{Renamed stations --- why the response lookup is epoch-aware}
KMA renames a station when it is re-installed: \code{BUS} $\rightarrow$ \code{BUS2} $\rightarrow$ \code{BUS3},
\code{DAG} $\rightarrow$ \code{DAG2}, \code{USN} $\rightarrow$ \code{USN2} --- \emph{identical coordinates},
strictly sequential response epochs. \code{responses.epoch\_families()} discovers these from the StationXML
coordinates (never from the trailing digit, so \code{GKP1} and \code{N001} are safe) and requires the epochs to
be \textbf{non-overlapping}, which correctly excludes co-located \emph{contemporaries} such as \code{CHS}/\code{CSO}
and \code{INC}/\code{INCA} --- genuinely different instruments that must never be aliased. 23 families, 49 codes.

\smallskip
This matters because of one archive quirk: \code{KS\_KG/BUS2/} holds \textbf{both} epochs' data, and its
miniSEED headers read \code{BUS} until 2015 --- long after the Nov-2010 rename. Association rewrites those
picks to \code{BUS2}, which is right for \emph{position} (identical coordinates, so location and HypoDD are
unaffected) but wrong for the \emph{response}, whose epochs are strict. Without the epoch-aware lookup, 89\% of
2010's BUS picks silently lost their response. An archive-wide header audit found \textbf{this is the only
station affected} --- \code{DAG2}, \code{ULJ2}, \code{USN2}, \code{GKP1/2}, \code{BUS3} and all GJ/NS dirs have
headers matching their directory.
\end{warnbox}

\begin{warnbox}{\code{ns\_derived/} rests on an assumption --- N2xx magnitudes are provisional}
The MARA dataless (2020-06-03) predates the \code{N201}--\code{N220} deployment (first data 2022.237), so those 20
stations have no authoritative response. All 600 authoritative NS RESP files are byte-identical apart from
station code and epoch start date --- one instrument configuration for the whole array (754.3 $\times$
$4\!\times\!10^{5} = 3.0172\!\times\!10^{8}$ counts/(m/s), exact) --- so \code{write\_derived\_ns} clones it for
the missing block, stamping each with that station's own first-day-on-disk as the start date. The assumption
is that the 2022 expansion kept the same configuration; the N2xx stations do record at the same 200\,Hz, in
3 components, at comparable count amplitudes, but there is no manufacturer record here to confirm it. Every
generated file carries a \code{DERIVED by ufpipe.responses} header. Delete the directory and re-derive the
moment real metadata arrives.
\end{warnbox}

%================================================================================
\section{How to run}

\begin{keybox}{Notebooks --- one tidy folder per year}
\code{python notebooks/make\_year\_notebooks.py -{}-years 2010-2024} builds, for each year,
\code{notebooks/<year>/00.Run\_pipeline\_<year>.ipynb} (the 6-stage cockpit: run + check + PyGMT maps, every
knob in the top parameters cell) and \code{01.Record\_sections\_<year>.ipynb} (record sections with associated
P/S picks + a dedicated augmented-events figure), with YEAR/MODEL pre-filled.\\[1mm]
\code{python notebooks/make\_year\_notebooks.py -{}-status} prints the per-year progress table
(picks / assoc / located / QC / dt.cc) read from disk --- the way to manage a multi-year run.
Kernel \code{base}; the emitted notebooks are gitignored (the builders are tracked).
\end{keybox}

\begin{keybox}{Standard runs (from the repo root)}
\begin{tabular}{@{}>{\ttfamily}p{8.4cm}>{\raggedright\arraybackslash}p{6.2cm}@{}}
python -m ufpipe.run\_pipeline -{}-model original -{}-years 2024 & one year, all 6 stages \\
python -m ufpipe.run\_pipeline -{}-model original -{}-years 2010-2024 & the whole record \\
python -m ufpipe.run\_pipeline -{}-model original -{}-years 2015,2016 -{}-velmodel kim1983 & specific years, kim1983 \\
python -m ufpipe.run\_pipeline -{}-model original -{}-years 2024 -{}-stage-from association & resume (picks already done) \\
python -m ufpipe.run\_pipeline -{}-model original -{}-years 2024 -{}-days 1-3 & detection on a few days (testing) \\
python -m ufpipe.run\_pipeline -{}-model original -{}-years 2024 -{}-strict & strict association gate \\
python -m ufpipe.run\_pipeline -{}-model original -{}-years 2016 -{}-stage-from relocate -{}-through dtcc & relocation only (dt.cc; needs feeder) \\
\end{tabular}
\end{keybox}

\noindent\textbf{CLI flags:} \code{-{}-model} · \code{-{}-years} (\code{'2010-2024'} or \code{'2015,2016'}) ·
\code{-{}-velmodel} · \code{-{}-days} · \code{-{}-device \{cuda,cpu\}} · \code{-{}-workers} · \code{-{}-strict} ·
\code{-{}-force} (allow writing to \code{stead}) ·\\ \code{-{}-stage-from
\{detection,association,augment,phs,locate,relocate\}} (start mid-chain) ·
\code{-{}-gate}/\code{-{}-pick-match-tol}/\dots\ (association, \S4.2) ·
\code{-{}-qc}/\code{-{}-xcorr} (relocation, \S7) · \code{-{}-no-regen-sta} (keep an existing station file).

\begin{warnbox}{Shared 64-core server — always pin cores}
Detection sizes its preprocessing pool and \code{torch} threads from the process CPU \emph{affinity}, capped by
\code{MAX\_CORES=24}. Launch under \code{taskset} so the whole job stays within budget:\\[1mm]
\hspace*{3mm}\code{OMP\_NUM\_THREADS=1 taskset -c 0-7 python -m ufpipe.run\_pipeline -{}-model original -{}-years 2010-2024}\\[1mm]
Without pinning, PhaseNet+ grabbed $\sim$49 cores / 193 threads and starved everything. Inference is
\textbf{GPU-preferred} (warns loudly, never silently falls back to CPU).
\end{warnbox}

\noindent Detection is \textbf{idempotent} — it skips station-days whose picks already exist, so runs resume
safely after an interruption.

\subsection{Running each stage separately}
Every \code{ufpipe} stage is also a standalone CLI (\code{-{}-model -{}-year}), so you can run detection alone,
then association alone, etc. Order is \code{detection $\rightarrow$ association $\rightarrow$ augment
$\rightarrow$ phs $\rightarrow$ locate}. Each reads the previous stage's output from disk.

\begin{center}
\small
\begin{tabular}{@{}>{\raggedright\arraybackslash}p{2.4cm}>{\ttfamily\raggedright\arraybackslash}p{12.0cm}@{}}
\toprule
Stage & \normalfont standalone command \\
\midrule
Detection & python -m ufpipe.detection -{}-model original -{}-year 2024 [-{}-networks KS,KG,GJ,NS] [-{}-days 1-5] [-{}-workers 8] \\
Association & python -m ufpipe.association -{}-model original -{}-year 2024 [-{}-networks KS,KG,GJ,NS] [-{}-workers 8] [-{}-strict] \\
Augment & \normalfont(runs inside \code{run\_pipeline}; or call \code{core.run\_augment\_year} from Python) \\
PHS & python -m ufpipe.make\_phs -{}-model original -{}-year 2024 \\
Location & python -m ufpipe.run\_hypoinverse -{}-model original -{}-year 2024 -{}-velmodel kim2011 \\
\bottomrule
\end{tabular}
\end{center}

Or drive the same stage functions from Python / a notebook:
\begin{quote}\small\ttfamily
from ufpipe import core, config\\
core.run\_detection\_year("original", 2024, days=range(1,6))\\
events, assign = core.run\_association\_year("original", 2024)\\
core.write\_phs("original", 2024); core.run\_hypoinverse\_year("original", 2024, velmodel="kim2011")
\end{quote}

\subsection{Relocation — running the last stage (stage 6)}
The relocate stage is \textbf{self-fed}: it needs only ufpipe's own per-year detection + association
(\code{pyocto\_kim2011\_<year>.csv} + assignment). It builds the reloc input store (event-idx SAC + station table)
from those via \code{ufpipe.reloc\_inputs}, then hands off to the driver \code{run\_picker\_reloc.py -{}-skip-build}
(scaffold $\rightarrow$ HYPOINVERSE $\rightarrow$ QC $\rightarrow$ rereference $\rightarrow$ GPU xcorr $\rightarrow$
HypoDD dt.cc). If a year's association is missing, the stage preflights and prints exactly what to run.

\begin{center}
\small
\begin{tabular}{@{}>{\raggedright\arraybackslash}p{2.9cm}>{\ttfamily\raggedright\arraybackslash}p{11.5cm}@{}}
\toprule
Step & \normalfont command \\
\midrule
Detection & conda run -n eqnet python -m ufpipe.detection -{}-model original -{}-year 2016 \\
Association & conda run -n base python -m ufpipe.association -{}-model original -{}-year 2016 \\
Reloc (to QC) & python -m ufpipe.run\_pipeline -{}-model original -{}-years 2016 -{}-stage-from relocate -{}-through hypoinverse \quad\normalfont(fast) \\
Reloc (full) & python -m ufpipe.run\_pipeline -{}-model original -{}-years 2016 -{}-stage-from relocate -{}-through dtcc \quad\normalfont(+ xcorr$\rightarrow$HypoDD dt.cc; $\sim$6\,h dense) \\
\bottomrule
\end{tabular}
\end{center}
\code{-{}-through hypoinverse} stops at the QC'd absolute-location subset; \code{-{}-through dtcc} continues to the
double-difference relocation. \code{-{}-clean-cache} deletes the $\sim$tens-of-GB interp cache after dt.cc
(essential for multi-year runs). Two override flags expose the stage-6 knobs per run (defaults = the
validated values): \code{-{}-qc "erh=5,erz=5,gap=270,num=5,rms=1.0"} (the \code{uf\_cluster.QC} gate) and
\code{-{}-xcorr "interp\_hz=1000,fmin=5,fmax=20,cc\_threshold=0.7,pre=0.5,post=0.5,margin=0.5"} (the dt.cc
cross-correlation). The HypoDD adaptive-damping/ISTART=2 scheme and the full-run-HYPOINVERSE injection are
fixed invariants (not flag-settable). The cockpit notebook (\code{notebooks/00.Run\_yearly\_pipeline.ipynb})
sets all of these from its parameters cell. \code{detection\_test/lib/} is DEPRECATED (see
\code{lib/DEPRECATED.md}); the relocation internals are documented in
\code{src/ufpipe/reloc\_driver/PIPELINE.md} + \code{study\_guide.pdf}.

%================================================================================
\section{Outputs \& layout}

All pipeline products go under \code{outputs/models/<model>/} (regenerable; gitignored):
\begin{center}
\small
\begin{tabular}{@{}ll@{}}
\toprule
Path & Contents \\
\midrule
\code{.../detection\_location/<year>/picks/} & per-station-day picks (stage 1) \\
\code{.../pyocto/pyocto\_\{vm\}\_\{year\}.csv} & associated events (stage 2) \\
\code{.../pyocto/pyocto\_assignment\_\{vm\}\_\{year\}.csv} & pick$\rightarrow$event assignment (authoritative) \\
\code{.../station\_table/stations\_\{year\}.csv} & per-year station roster (discovered from metadata) \\
\code{.../HypoInv/PHS/UF\{year\}.phs} & HYPOINVERSE phase file (stage 4) \\
\code{.../HypoInv/\{vm\}/UF\{year\}.\{sum,arc,prt\}} & located catalog + errors (stage 5) \\
\code{outputs/reloc/reloc\_\{year\}\_uf[\_\{model\}]/results/} & dt.cc-relocated catalog + links (stage 6) \\
\bottomrule
\end{tabular}
\end{center}
\noindent \code{\{vm\}} = \code{config.PYOCTO\_VELMODEL} = \textbf{kim2011} for current daily-chunked runs;
legacy pre-2026-07 whole-year artifacts on disk are named \code{kim1983}.

\medskip
\noindent\textbf{Inputs} (metadata by kind, 2026-07): waveforms in \code{KS\_KG/ GJ/ NS/ NS\_100hz/} (station
dirs at each root); station tables in \code{data/metadata/stations/\{ks\_kg,gj,ns,kigam\}/} (+ the per-year
derived cache \code{stations/derived/}); instrument responses for all four networks in
\code{data/metadata/responses/\{master,fetched,gj,ns,ns\_derived\}/} (\S4.4); velocity
model in \code{data/metadata/velocity/}; HYPOINVERSE control inputs (STA/\code{kim*}) in \code{data/hypoinv/}.

%================================================================================
\section{Stage 6: relocation — HypoDD dt.ct + dt.cc}

The 6th ufpipe stage: turn the QC'd absolute-location catalog into a \textbf{double-difference relocated} catalog
for sub-100\,m relative precision.

\begin{quote}\ttfamily\small
python -m ufpipe.run\_pipeline -{}-model original -{}-years 2016 -{}-stage-from relocate -{}-through dtcc -{}-clean-cache
\end{quote}

\code{ufpipe.relocate} (\code{src/ufpipe/relocate.py}) is \emph{self-fed}: it preflights ufpipe's OWN per-year
association, builds the reloc inputs from it (\code{ufpipe.reloc\_inputs} --- SAC store, pyocto tables, station
table, merged archive), then invokes the validated orchestrator
\code{src/ufpipe/reloc\_driver/run\_picker\_reloc.py -{}-skip-build}, which drives the external
HypoDD/xcorr engine. Internal stages (all KS/KG/GJ/NS):

\begin{center}
\small
\begin{tabular}{@{}c>{\raggedright\arraybackslash}p{3.0cm}>{\raggedright\arraybackslash}p{9.2cm}@{}}
\toprule
\# & Stage & What it does \\
\midrule
1 & SAC store & event-idx-keyed native-rate SAC waveforms from the associated picks (\code{event\_sac/}) \\
2 & HYPOINVERSE & absolute location of all UF-box events (kim2011) $\rightarrow$ \code{.sum} \\
3 & QC & \code{uf\_cluster.QC}: keep \code{erh$<$5, erz$<$5, gap$<$270, num$>$5, rms$<$1.0} \\
4 & inject full \code{.sum}/\code{.arc}/STA & QC subset reuses the full-run HYPOINVERSE solution + station file (no redundant re-run) \\
5 & ph2dt / dtct & build \code{event.dat} + catalog differential times \code{dt.ct} from the injected \code{.arc} \\
6 & re-reference & stamp each SAC's origin from the (full-run) HYPOINVERSE \code{.sum} \\
7 & xcorr (GPU) & cross-correlation differential times \code{dt.cc} (band 5--20 Hz, interp 1000 Hz, cc$\ge$0.7) \\
8 & HypoDD & kim2011, ISTART=2, ISOLV=2 (LSQR), adaptive damping (CND 60--80), HypoDD v2.1beta (MAXDATA=15M) $\rightarrow$ \code{hypoDD.reloc} \\
\bottomrule
\end{tabular}
\end{center}

\begin{keybox}{Relocation outputs (symlinked into the working dir)}
Each run publishes \code{outputs/reloc/reloc\_<year>\_uf[\_<p>]/results/} with symlinks to the external run tree:
\code{hypoDD.reloc.dtcc} (dt.ct+dt.cc — \textbf{the deliverable}), \code{HypoInv.full.sum}, \code{HypoInv.qc.sum}, \dots
``dt.cc-resolved'' = events keeping $\ge$1 surviving cross-correlation link — the highest-precision subset.
\code{hypoDD.reloc.dtct} (dt.ct-only) is an \emph{optional diagnostic baseline}: it re-runs HypoDD on catalog
differential times only (heavier, no cc) and is non-fatal + time-bounded, so it may be absent when slow — the
dt.cc result is unaffected.
\end{keybox}

\begin{warnbox}{The invariant that matters (see PIPELINE.md)}
Origins flow into dt.cc: \code{rereference} stamps SAC origins from a \code{.sum}, and xcorr stores
$dt=(t_1{+}\text{shift}{-}ot_1)-(t_2{-}ot_2)$ — subtracting those origins. Use the \emph{full-run} HYPOINVERSE
solution QC gated on (never a redundant re-run), or the dt.cc \emph{values} are corrupted. This is enforced by
\code{run\_picker\_reloc.py::inject\_full\_hypoinverse}.
\end{warnbox}

%================================================================================
\section{Robustness \& invariants}

\begin{itemize}[leftmargin=6mm,itemsep=1mm]
\item \textbf{Network growth is automatic.} Stations are \emph{discovered from StationXML + on-disk coverage}
each month; the deployed network for that period is what enters the run. No hard-coded station lists.
\item \textbf{Mixed sampling rates handled.} $>$100 Hz stations (KG/NS at 200 Hz) are anti-alias-lowpassed then
resampled to the pickers' 100 Hz. The anti-alias filter is applied BEFORE down-sampling (never after).
\item \textbf{Velocity model, thresholds, and gates are epoch-invariant by choice} $\Rightarrow$ the catalog stays
internally comparable across years.
\item \textbf{Idempotent \& resumable} at every stage; \code{-{}-stage-from} restarts mid-chain.
\item \textbf{Do not edit the \code{stead} reference run} (it is the baseline). Scripts refuse \code{-{}-model
stead} writes unless \code{-{}-force}.
\end{itemize}

%================================================================================
\section{Troubleshooting / provenance checks}

\begin{center}
\small
\begin{tabular}{@{}>{\raggedright\arraybackslash}p{6.0cm}>{\raggedright\arraybackslash}p{8.5cm}@{}}
\toprule
Symptom & Check \\
\midrule
Detection finds no waveforms & \code{python -c "import ufpipe.stations as s; print(s.build\_year\_table(YEAR).net.value\_counts())"} --- the per-year multi-network table must list stations, and each row's \code{archive} dir must exist. \\
\code{ImportError: ufpipe} / \code{uflib} & \code{pip install -e . -{}-no-deps} in \textbf{both} \code{eqnet} and \code{base}; \code{import ufpipe.config; print(\_\_file\_\_)} must be \code{src/ufpipe}. \\
\code{ModuleNotFoundError: pyocto} & you ran association in \code{eqnet} --- PyOcto lives in \code{base}. Detection $\to$ \code{eqnet}; association $\to$ \code{base}. \\
\code{InternalMSEEDWarning: ...Steim2 failed} & corrupt Steim2 record(s) in the archive (e.g.\ KG.MKL 2010): the samples still decode; the reader collapses the per-record spam into one summary line per station-day and reads anyway. Informational --- the named station-days are the QC signal. \\
GJ/NS station silently yields no picks & mid-day sampling-rate changes (100/200/1000\,Hz) and int32/float64 records refuse to merge --- handled since 2026-07 (rate + dtype unified \emph{before} merge in both readers). \\
\code{relocate} aborts at preflight & ufpipe's own per-year association is missing for that (model, year) --- run \code{ufpipe.detection} (eqnet) then \code{ufpipe.association} (base); the preflight prints the exact commands. \\
Wrong \code{pipeline} imported & \code{ufpipe} is renamed precisely to avoid the collision with the external korea-cluster-relocation \code{pipeline}. Verify \code{ufpipe.config.\_\_file\_\_} is under \code{src/}. \\
Events look over/under-associated & \code{ASSOC\_GATE} vs \code{-{}-strict} (\code{ASSOC\_GATE\_STRICT}); \code{ASSOC\_PICK\_MATCH\_TOL} is the primary knob. \\
GPU not used & inference is GPU-preferred and warns loudly; check the torch/CUDA install (never silently CPU). \\
Sparse early years, few events & expected --- KS/KG-only network (2010--2015) has larger azimuthal gaps; GJ joins 2016, the dense NS array 2017+; the QC gate honestly reflects what the network resolves. \\
\code{.sum} stops mid-year; months missing & \code{hyp1.40} \textbf{segfaulted} partway (2016 once died on 25 Jun, dropping the whole Gyeongju sequence): it crashes in \code{hytrl\_} on any event whose P picks are ALL weight code 4 (only possible under the probability weight scheme, every P below the last \code{PHS\_WEIGHT\_BINS} threshold). Fixed 2026-07: \code{write\_phs} skips such unlocatable events (logged to \code{PHS/UF<year>.phs.skipped}) and \code{run\_hypoinverse\_year} now raises on a nonzero exit code instead of shipping the partial \code{.sum}. If you see this, regenerate from stage \code{phs}. \\
\bottomrule
\end{tabular}
\end{center}

\vspace{2mm}\hrule\vspace{2mm}
{\small\color{soft} Companion docs: \code{docs/overview.md}, \code{docs/pipeline.md}, \code{docs/how-to-run.md},
\code{docs/directory-structure.md} · \code{src/ufpipe/README.md} · relocation pipeline:
\code{src/ufpipe/reloc\_driver/PIPELINE.md} + \code{study\_guide.pdf}.}

\end{document}
